\documentclass[10pt,a4paper]{amsart}
\usepackage[margin=19mm]{geometry}
\usepackage{amsmath,amssymb,mathtools}
\usepackage[scr=rsfs,cal=euler]{mathalfa}
\usepackage{xfrac}
\usepackage[unicode,hidelinks,bookmarks=false]{hyperref}
\newcommand{\ie}{i.e.\ }
\newcommand{\cf}{cf.\ }
\newcommand{\resp}{respectively\ }
\newcommand{\Subsection}{\subsection}
\providecommand{\nobreakdash}{\nobreak\hbox{-}}
\setlength{\emergencystretch}{2em}
\multlinegap0pt
\usepackage{enumerate}
\usepackage{amssymb}
\usepackage{tikz}
\newtheorem{theorem}{Theorem}
\newtheorem{lemma}{Lemma}
\title{A note on the $k$-defect number: \\ Vertex Coloring with a Fixed Number of Monochromatic Edges}
\author{Eunice Mphako-Banda, Christo Kriel, Alex Alochukwu}
\date{Source: arXiv:2510.00712v2 (CC BY 4.0)}
\begin{document}
\maketitle

\noindent\emph{Source and licence.} A note on the $k$-defect number: \\ Vertex Coloring with a Fixed Number of Monochromatic Edges. Version arXiv:2510.00712v2 (CC BY 4.0), distributed by arXiv under the Creative Commons Attribution 4.0 International licence, http://creativecommons.org/licenses/by/4.0/, as recorded in the arXiv licence metadata for this exact version and shown on the abstract page https://arxiv.org/abs/2510.00712v2. Full licence notice: \texttt{licenses/arxiv-2510.00712-CC-BY-4.0.txt}.\par\smallskip

\noindent\textit{Editorial scope.} Counted unit: the article's theorem th:kdefectcontr (source0.tex lines 375--377). The article's complete original proof of it is delivered with it (source0.tex lines 378--394, one \texttt{proof} environment of 17 lines). No mathematics is written, repaired or re-derived: the statement and the proof are the article's own lines, in the article's own notation, and no retained line is substituted or re-flowed.  Retained uncounted as the source-local context the proof needs: lines 241--245; lines 370--374.  Nothing was omitted from any retained span, so the package carries no omission record.  No retained line contains a citation, so the wrapper carries no bibliography and no bibliography entry.  No retained line contains a figure environment, an image-inclusion command or drawing code, so no image is delivered and the package declares no asset dependency.  Editorial formatting only, in addition: an AMS \texttt{amsart} preamble that restates the article's own declarations and macros used by the retained lines, namely the environments \texttt{theorem}, \texttt{lemma} on separate counters, together with the standard packages those lines need.  The retained environments print in this document as Theorem 1, Lemma 1 in order of appearance, whereas the article prints the same items with the numbering of its own full document.  The complete native source of the article as distributed in the arXiv e-print for this exact version is bundled unchanged as \texttt{sources/186152/source0.tex}.
\begin{theorem}
\label{prop2}
Let $G$ be a graph. Then  $$\phi_{k}(G;\lambda) = \sum_{X \in L(G),\vert X \vert = k} \chi (G/X; \lambda)$$
if $G$ has at least one closed set  of size $k$. Otherwise $\phi_{k}(G;\lambda) = 0.$
\end{theorem}

\begin{lemma}
\label{lem1}
 The chromatic polynomial of $G$ can be factorised as 
 $\chi (G; \lambda)= \lambda (\lambda -1)(\lambda -2) \cdots (\lambda -r)[Q(\lambda)]$ such that $\chi(G)=r+1$  and $Q(\lambda)$ is a polynomial without positive integer roots.   
\end{lemma}

\begin{theorem}The $k$-defect number $\phi_k(G) = \chi(G/X)$, where $\chi(G/X)$ is the  minimum chromatic number  over all minors of $G$ obtained by contracting closed sets of size $k.$
\label{th:kdefectcontr}
\end{theorem}

\begin{proof} 
By Theorem~\ref{prop2} we have,
$\displaystyle \phi_{k}(G;\lambda) = \sum_{X\in L(G),\vert X \vert = k} \chi (G/X; \lambda)$ if $G$ has at least one closed set  of size $k.$ Thus,  if $L(G) = \{X_1, X_2, \cdots , X_n\}$ then we can rewrite 
\begin{eqnarray*}
   \phi_{k}(G;\lambda) &=& \sum_{X_i} \chi(G/X_i; \lambda).\\
   &=&  \sum\lambda(\lambda -1)(\lambda -2) \cdots (\lambda -r_i)  Q_i(\lambda) \text{ by Lemma~\ref{lem1} }.\\
\end{eqnarray*}

Hence, we can rewrite
\begin{eqnarray*}
   \phi_{k}(G;\lambda) 
   &=& \lambda(\lambda -1)(\lambda -2) \cdots (\lambda -r) T_i(\lambda) Q_i(\lambda)
\end{eqnarray*}
where $r$ is the smallest $r_i$ and $T_i(\lambda) = (\lambda -r_i+1) \cdots (\lambda -r_i+t)$ for some $t \geq 0.$

    
\end{proof}

\end{document}
